\documentclass[11pt]{article}

% ===============================
% Geometry and Layout
% ===============================
\usepackage[margin=1.1in]{geometry}
\usepackage{setspace}
\setstretch{1.15}

% ===============================
% Engine and Fonts (LuaLaTeX)
% ===============================
\usepackage{fontspec}
\setmainfont{Libertinus Serif}
\setsansfont{Libertinus Sans}
\setmonofont{Libertinus Mono}

% ===============================
% Mathematics and Formalism
% ===============================
\usepackage{amsmath,amssymb,amsthm,mathtools}

\newtheorem{definition}{Definition}
\newtheorem{proposition}{Proposition}
\newtheorem{theorem}{Theorem}
\newtheorem{lemma}{Lemma}

% ===============================
% Microtypography and Tables
% ===============================
\usepackage{microtype}
\usepackage{booktabs}

% ===============================
% Hyperlinks
% ===============================
\usepackage[hidelinks]{hyperref}


% ===============================
% Title Information
% ===============================
% ===============================
% Title Information
% ===============================
\title{Coherence Before Engagement:\\
An Architecture for Persistent Meaning in Digital Systems}

\author{Flyxion}

\date{\today}

\begin{document}
\maketitle

\begin{abstract}
Contemporary social media platforms have become foundational infrastructure for public
discourse, knowledge exchange, and collective sense--making. Despite their scale and
pervasiveness, these systems exhibit persistent and well--documented pathologies:
semantic drift, performative identity, misinformation amplification, and the erosion of
interpersonal trust. This paper argues that these failures are not the result of deficient
moderation policies or misaligned recommendation algorithms, but instead arise from
fundamental architectural commitments embedded in feed--based, engagement--optimized
platforms.

We present a coherence--first alternative to social media grounded in a formal semantic
substrate. In this architecture, the primary unit of interaction is not a post or message,
but a structured semantic object with explicit modality requirements, provenance, and
bounded entropy. Identity is defined by persistence and traceable transformation history
rather than visibility or popularity metrics. Evolution of meaning proceeds through
typed, entropy--bounded transformations that preserve coherence over time.

Rather than proposing incremental reforms to existing platforms, this work outlines a
categorically different model of social computation. We formalize core concepts, analyze
their information--theoretic properties, provide concrete instantiations, and examine
scalability, governance, and implementation considerations. The result is a principled
framework for large--scale digital discourse in which trust and understanding emerge as
stable structural properties rather than ephemeral signals.
\end{abstract}

\subsection{Contributions}

This paper makes the following contributions. First, it provides a formal critique of
feed--based social media architectures grounded in information theory and knowledge
representation. Second, it introduces a coherence--first semantic substrate in which
meaning is treated as a structured, evolving object rather than a sequence of messages.
Third, it formalizes entropy--bounded semantic transformation and provenance--anchored
identity, and analyzes their theoretical properties. Fourth, it presents concrete
instantiations and worked examples demonstrating how discourse, revision, and discovery
operate without feeds or engagement metrics. Finally, it examines scalability,
governance, and implementation considerations, identifying both strengths and open
challenges of the proposed approach.

\subsection{Roadmap}

The remainder of this paper is organized as follows. Section~2 reviews related work,
situating the proposed architecture relative to existing social platforms, collaborative
systems, and academic literature in CSCW, HCI, and semantic knowledge systems.
Section~3 develops the theoretical foundations of coherence--first design, including
formal definitions of semantic objects, entropy--bounded evolution, and provenance
graphs. Section~4 presents concrete instantiations and worked examples illustrating how
these concepts operate in practice. Section~5 addresses challenges and limitations,
including scalability, onboarding, and adversarial behavior. Section~6 formalizes key
claims and provides quantitative analyses. Section~7 examines the constrained role of
artificial intelligence. Section~8 elaborates governance mechanisms and invariants.
Section~9 discusses implementation considerations. Section~10 concludes with a summary
of contributions and a research agenda for future work.

\subsection{A Motivating Example}

Consider a long--running public discussion concerning a scientific or policy question.
On contemporary social media, this discussion fragments into thousands of posts,
responses, and quote--reposts, each optimized for visibility at the moment of publication.
Corrections are rarely propagated with the same reach as initial claims, and revisions
are indistinguishable from contradictions. Participants are rewarded for novelty and
provocation rather than sustained engagement with the evolving substance of the topic.

In contrast, a coherence--first system represents the discussion as a small number of
persistent semantic objects. Claims, evidence, counterarguments, and revisions are
attached to these objects through typed transformations with explicit provenance. The
discussion evolves by accumulation and refinement rather than displacement. Visibility
is decoupled from immediacy, and trust accrues through demonstrated consistency over
time.

The remainder of this paper develops the formal and practical implications of adopting
such an architecture.

\section{Related Work}

The proposal advanced in this paper intersects multiple lines of prior work, including
commercial social media platforms, collaborative knowledge systems, federated networks,
and academic research in computer--supported cooperative work (CSCW), human--computer
interaction (HCI), semantic web technologies, and knowledge representation. This section
situates the coherence--first approach relative to these domains and clarifies how it
differs from both existing platforms and incremental reform efforts.

\subsection{Feed--Based Social Media Platforms}

Dominant social media platforms such as Twitter/X, Reddit, and Facebook share a common
architectural pattern: content is produced as discrete posts and delivered to users via
feeds ordered by recency, engagement, or learned relevance. Identity is represented by
profiles augmented with follower counts, karma, or verification markers. Visibility and
influence are mediated by algorithmic ranking systems that prioritize interaction volume
over semantic structure.

Extensive empirical work has documented the consequences of these design choices. Studies
in CSCW and computational social science have shown that engagement--based ranking
amplifies emotionally charged content, accelerates polarization, and privileges novelty
over correction or synthesis. Once introduced, misinformation persists even after
retraction, as corrective posts rarely propagate through the same network paths as the
original claims. These phenomena are not anomalies but stable equilibria of
attention--optimized systems.

Incremental reforms within this paradigm have focused on moderation, labeling, or
algorithmic adjustment. While such interventions may reduce specific harms, they do not
alter the underlying dynamics produced by feed ordering and popularity metrics. The
coherence--first approach diverges by rejecting the feed as a primary organizing
structure and eliminating engagement metrics entirely, rather than attempting to refine
their use.

\subsection{Community Aggregation and Forum Systems}

Platforms such as Reddit and Stack Exchange introduce additional structure through topic
segmentation, threaded discussion, and reputation systems. These designs mitigate some
issues associated with unstructured feeds by localizing discourse and enabling deeper
engagement within bounded communities. However, they continue to rely on post--level
interaction and popularity signals, such as upvotes, karma, and accepted answers.

Research on online forums highlights tradeoffs inherent in these systems. Reputation
mechanisms can incentivize expertise and sustained participation, but they also encourage
strategic behavior and gatekeeping. Threads remain ephemeral, with discussions reset as
new posts displace older ones. Knowledge accumulates unevenly and often requires manual
curation or external summarization to remain accessible.

The coherence--first model differs by treating discussion topics as persistent semantic
objects rather than ephemeral threads. Contributions do not compete for visibility within
a feed or forum, but attach directly to evolving objects whose structure constrains how
discourse unfolds.

\subsection{Federated and Decentralized Networks}

Federated platforms such as Mastodon and Bluesky aim to address governance and moderation
concerns by decentralizing control. These systems distribute authority across instances
or protocols while preserving familiar social media interactions, including posts,
feeds, and follower relationships.

Decentralization addresses important issues of platform control and censorship but leaves
the semantic architecture largely unchanged. Feed--based delivery, engagement signals,
and performative identity remain central. As a result, many of the same coherence and
trust problems observed in centralized platforms reappear in federated settings.

The coherence--first approach is orthogonal to decentralization. While compatible with
distributed deployment, it targets the semantic and computational substrate rather than
the ownership or governance layer. Decentralization without semantic constraints risks
replicating existing failure modes at smaller scales.

\subsection{Newsletter and Long--Form Publishing Platforms}

Platforms such as Substack emphasize long--form content and subscription--based
distribution, offering an alternative to short--form social media. By decoupling revenue
from engagement metrics and emphasizing author--reader relationships, these systems
encourage slower, more deliberate communication.

However, long--form publishing platforms retain a broadcast--oriented model. Content is
still produced as discrete artifacts, revisions are limited, and discourse occurs largely
through comments or external channels. Semantic relationships between texts are implicit
rather than formally represented.

The coherence--first model generalizes beyond long--form publishing by supporting
structured evolution of ideas across modalities and time, rather than privileging a
single medium or authorial voice.

\subsection{Collaborative Knowledge Systems}

Wikis represent an early attempt to support cumulative knowledge construction. By
allowing collaborative editing of shared documents, wikis reduce fragmentation and
encourage synthesis. Version histories provide a limited form of provenance, enabling
reversion and accountability.

Despite these strengths, wikis impose a monolithic document structure that struggles to
accommodate disagreement, parallel perspectives, and ongoing revision. Semantic
relationships between pages are informal, and conflicts are resolved socially rather than
mechanically. Empirical studies of wiki governance document persistent tensions between
openness and control.

Version control systems such as Git provide more rigorous provenance tracking and support
branching and merging of parallel developments. These systems excel at managing structured
artifacts such as source code but are ill--suited to general discourse. Line--based diffs
and merge conflicts capture syntactic changes rather than semantic ones, and non--textual
modalities are poorly supported.

The coherence--first approach draws inspiration from both wikis and version control while
addressing their limitations. It generalizes provenance and branching to semantic objects
across modalities and introduces entropy--bounded merging to preserve meaning rather than
syntax.

\subsection{Semantic Web and Knowledge Representation}

The semantic web and knowledge representation communities have long sought to formalize
meaning through ontologies, graphs, and logical inference. Technologies such as RDF, OWL,
and knowledge graphs provide expressive frameworks for representing structured information
and relationships.

While powerful, these systems typically require significant upfront formalization and
expertise. They are optimized for machine reasoning rather than human discourse and lack
mechanisms for incremental, conversational evolution of meaning. As a result, they have
seen limited adoption as substrates for everyday communication.

The coherence--first architecture occupies a middle ground. It adopts formal structure
where necessary to preserve coherence but does not require full logical formalization of
all content. Semantic constraints operate locally and incrementally, allowing discourse to
remain flexible while preventing unbounded drift.

\subsection{Architectural Alternatives versus Incremental Reform}

A recurring theme in prior work is the attempt to repair social media pathologies through
incremental changes: improved moderation, transparency tools, algorithmic tweaks, or
decentralized governance. While valuable, these efforts accept the core primitives of
posts, feeds, and engagement metrics as given.

This paper advances a different stance. It argues that these primitives are the source of
systemic failure and must be replaced rather than refined. The coherence--first approach
constitutes an architectural alternative, redefining the units of communication, identity,
and relevance at the substrate level.

In doing so, it aligns with a growing body of research emphasizing the importance of
structural constraints in shaping collective behavior. Rather than optimizing outcomes
within a fixed architecture, it seeks to redesign the architecture itself to make certain
failure modes structurally impossible.

\section{Theoretical Foundations}

This section develops the formal foundations of a coherence--first semantic substrate.
We introduce semantic objects as the primary units of meaning, define entropy--bounded
evolution as a constraint on semantic transformation, specify modality requirements, and
formalize provenance as a graph--structured invariant. Together, these elements define a
computational framework in which meaning can evolve without unbounded drift.

\subsection{Semantic Objects}

We begin by defining the fundamental unit of representation.

\begin{definition}[Semantic Object]
A \emph{semantic object} is a structured entity
\[
\sigma = (I, \mathcal{M}, C, P, E)
\]
where $I$ is an immutable identity, $\mathcal{M}$ is a finite set of required modalities,
$C$ is a partial function assigning content to modalities, $P$ is a provenance graph
recording the derivation history of $\sigma$, and $E$ is a non--negative real number
representing the semantic entropy of the object.
\end{definition}

Intuitively, a semantic object represents a unit of meaning that may span multiple
modalities, such as text, data, code, audio, or formal proof. Unlike a post or message, a
semantic object persists across time and accumulates structure through controlled
transformation.

\begin{definition}[Well--Formed Semantic Object]
A semantic object $\sigma$ is \emph{well--formed} if and only if for all $m \in
\mathcal{M}$, the content assignment $C(m)$ is defined.
\end{definition}

Well--formedness enforces modality completeness. For example, a claim that requires both
textual explanation and empirical evidence cannot be considered complete until both
modalities are populated.

\subsection{Examples and Counterexamples}

As a motivating example, consider a scientific hypothesis represented as a semantic
object with modalities $\{\text{text}, \text{data}, \text{analysis}\}$. A draft
hypothesis with only textual content is not well--formed. It becomes well--formed only
once supporting data and analysis are attached.

In contrast, a short social media post asserting the hypothesis without evidence would
constitute a semantic object with $\mathcal{M} = \{\text{text}\}$, thereby encoding a
different and weaker claim. This distinction prevents implicit inflation of meaning and
clarifies the scope of assertions.

\subsection{Semantic Entropy}

We now formalize the notion of semantic entropy.

\begin{definition}[Semantic Entropy]
Let $\sigma$ be a semantic object. The semantic entropy $E(\sigma)$ is a scalar quantity
that measures internal inconsistency, ambiguity, or unresolved divergence within the
object.
\end{definition}

Semantic entropy is not identical to Shannon entropy, though it may incorporate
information--theoretic components. Operationally, it bounds the degree to which a
semantic object can tolerate conflicting content, incomplete modalities, or unresolved
branches.

\begin{definition}[Entropy--Bounded Transformation]
A transformation $T$ applied to a semantic object $\sigma$ is \emph{entropy--bounded} if
\[
E(T(\sigma)) \leq E(\sigma) + \varepsilon
\]
for a fixed budget $\varepsilon \geq 0$ associated with $T$.
\end{definition}

Entropy budgets constrain how much disorder a transformation may introduce. Transformations
that exceed their budgets are rejected or deferred pending reconciliation.

\subsection{Semantic Transformations}

Transformations operate on semantic objects to produce new versions.

\begin{definition}[Semantic Transformation]
A semantic transformation is a partial function
\[
T : \Sigma \rightarrow \Sigma
\]
defined on the space of semantic objects $\Sigma$, subject to modality preservation and
entropy constraints.
\end{definition}

Examples of transformations include revision, extension, summarization, translation, and
formalization. Each transformation specifies its input and output modality requirements
and an associated entropy budget.

\begin{proposition}[Monotonic Provenance]
Every valid transformation $T(\sigma)$ extends the provenance graph $P$ of $\sigma$ by
adding a new node corresponding to $T$.
\end{proposition}

\begin{proof}
By construction, transformations are required to record their application as edges in
the provenance graph. Since identities are immutable and transformations are append--only,
the graph grows monotonically.
\end{proof}

\subsection{Provenance Graph Structure}

Provenance is represented as a directed acyclic graph.

\begin{definition}[Provenance Graph]
The provenance graph $P$ of a semantic object $\sigma$ is a directed acyclic graph whose
nodes correspond to transformation events and whose edges encode dependency relations
between transformations.
\end{definition}

The acyclicity of $P$ ensures that semantic evolution is temporally ordered and that
every state has a finite justification history.

\begin{theorem}[Traceability]
For any semantic object $\sigma$, every content element $C(m)$ admits a finite derivation
trace in $P$.
\end{theorem}

\begin{proof}[Sketch]
Since $P$ is acyclic and transformations are finite in number for any realized object,
each content element corresponds to a path from an initial state to the current state.
\end{proof}

\subsection{Coherence Enforcement}

Coherence is enforced mechanically through validation checks.

\begin{definition}[Coherent State]
A semantic object $\sigma$ is \emph{coherent} if it is well--formed and its semantic
entropy does not exceed a system--defined threshold $E_{\max}$.
\end{definition}

\begin{proposition}[Rejection of Unbounded Drift]
Unbounded semantic drift is structurally impossible in a system enforcing entropy--bounded
transformations.
\end{proposition}

\begin{proof}[Sketch]
Suppose an unbounded drift occurs. Then $E(\sigma)$ must increase without bound along a
sequence of transformations. This contradicts the existence of fixed entropy budgets and
the enforcement of $E_{\max}$.
\end{proof}

\subsection{Identity and Persistence}

Identity is defined at the level of semantic objects.

\begin{definition}[Persistent Identity]
The identity $I$ of a semantic object is immutable and preserved across all valid
transformations.
\end{definition}

Deletion is therefore interpreted not as erasure but as a transformation that marks an
object as inactive while preserving its provenance.

\begin{lemma}[Non--Repudiation]
No valid transformation can remove or alter prior provenance.
\end{lemma}

\begin{proof}
Transformations are append--only with respect to provenance. Any operation that attempts
to remove prior nodes violates the definition of a semantic transformation.
\end{proof}

This completes the formal foundation of coherence--first semantic evolution.

\section{Concrete Instantiation and Worked Examples}

To clarify how a coherence--first semantic substrate operates in practice, this section
presents a series of worked examples. These examples illustrate the lifecycle of semantic
objects, the application of entropy--bounded transformations, identity persistence under
revision and deletion, and navigation through semantic neighborhoods in contrast to
feed--based browsing.

\subsection{Lifecycle of a Semantic Object}

Consider an initial semantic object $\sigma_0$ representing a factual claim. Let the
required modality set be
\[
\mathcal{M} = \{\text{text}, \text{evidence}\}.
\]
At time $t_0$, only the textual modality is populated. The object is therefore not
well--formed.

A transformation $T_1$ attaches empirical evidence, yielding $\sigma_1$ such that
$C_{\sigma_1}(\text{evidence})$ is defined. Provided that $T_1$ satisfies its entropy
budget, $\sigma_1$ becomes well--formed and eligible for further transformation. The
provenance graph records $T_1$ as a child of $\sigma_0$.

A subsequent transformation $T_2$ revises the textual explanation for clarity. Because
$T_2$ preserves modality completeness and introduces only bounded entropy, the resulting
object $\sigma_2$ remains coherent. Importantly, $\sigma_0$, $\sigma_1$, and $\sigma_2$
share a common identity and are linked by provenance rather than replacing one another.

\subsection{Entropy--Bounded Revision}

Revision is a common operation in discourse and often a source of semantic instability in
feed--based systems. In the coherence--first model, revision is formalized as a
transformation with a tightly constrained entropy budget.

Let $T_r$ be a revision transformation with budget $\varepsilon_r$. Applying $T_r$ to
$\sigma$ yields $\sigma'$ only if
\[
E(\sigma') \leq E(\sigma) + \varepsilon_r.
\]
If a proposed revision introduces ambiguity, contradiction, or modality imbalance beyond
this threshold, it is rejected or flagged for mediation.

This mechanism distinguishes legitimate refinement from destabilizing alteration.
Incremental clarification typically falls within budget, while radical reinterpretation
requires explicit branching, preserving traceability.

\subsection{Branching and Conflict Resolution}

Conflicts arise when incompatible transformations are proposed. In such cases, the
provenance graph branches.

Let $\sigma$ admit two transformations $T_a$ and $T_b$ that cannot be reconciled within a
shared entropy budget. The system produces two descendant objects $\sigma_a$ and
$\sigma_b$, each inheriting the identity of $\sigma$ but diverging in provenance.

Rather than forcing premature resolution, the system allows both branches to persist.
Subsequent transformations may reconcile the branches by introducing additional evidence,
clarification, or scope restriction. Any reconciliation must itself be entropy--bounded
and is recorded as a transformation whose provenance references both branches.

This approach contrasts with feed--based systems, where conflicts are resolved implicitly
through visibility dynamics rather than explicit semantic structure.

\subsection{Deletion and Retraction}

Deletion presents a critical edge case. In social media, deletion often erases context,
undermining accountability and enabling revisionist narratives.

In a coherence--first system, deletion is modeled as a transformation $T_d$ that marks a
semantic object as inactive. The content remains accessible for provenance and reference,
but the object is no longer eligible for extension or endorsement.

Retraction is handled similarly but includes an explicit statement of invalidation. A
retraction transformation increases semantic entropy temporarily while introducing
constraints that prevent further use of the object as a premise. This preserves historical
traceability while signaling reduced validity.

\subsection{Identity Persistence Across Time}

Identity persistence follows directly from the immutability of semantic object identities.
All transformations preserve identity, ensuring that long--term evolution remains
traceable even as content changes.

This property supports longitudinal trust. Observers can inspect how an object has
changed, which transformations were applied, and which branches were reconciled or
abandoned. Trust is therefore grounded in demonstrated behavior over time rather than
static declarations or popularity signals.

\subsection{Semantic Neighborhood Navigation}

Discovery in a coherence--first system operates through semantic neighborhoods rather than
feeds. Users encounter semantic objects by traversing relationships defined by shared
modalities, transformation types, or provenance connections.

Navigation proceeds through deliberate exploration rather than passive consumption.
Because objects persist and accumulate structure, older but coherent material remains
discoverable without being displaced by newer content.

This mode of interaction privileges understanding over immediacy and reduces the
attention fragmentation characteristic of infinite scroll interfaces.

\subsection{Comparison with Feed Browsing}

In feed--based systems, relevance decays rapidly with time, and visibility is governed by
engagement metrics. In contrast, semantic neighborhoods decouple relevance from recency.
Objects remain accessible as long as they remain coherent.

This shift alters user incentives. Contributions that improve coherence, clarify meaning,
or reconcile conflict are rewarded structurally, while reactive or sensational content
gains no inherent advantage.

These examples demonstrate that the coherence--first model is not merely theoretical but
operationally viable, offering concrete mechanisms for discourse, revision, and discovery
without reliance on feeds or engagement metrics.

\section{Challenges, Limitations, and Structural Tradeoffs}

The coherence--first semantic substrate described in this paper is not merely an alternative
content organization scheme, but a reconfiguration of social interaction around spatial,
computational, and educational primitives. In this section, we revisit challenges and
limitations under the assumption that the platform functions simultaneously as a
programming environment, a persistent semantic memory, and a navigable social space
embedded in a high--dimensional latent geometry.

\subsection{Scalability in a Spatial Semantic Medium}

Concerns regarding scalability are often framed in terms of throughput and user count,
implicitly assuming feed--based dissemination as the baseline. In a spatial semantic
system, scale is expressed differently. Growth manifests as expansion of a latent semantic
manifold rather than increased velocity within a global feed.

Knowledge objects are not broadcast to all users. Instead, they occupy stable locations in
semantic space. As the system grows, density increases locally rather than globally, and
users encounter content by navigating regions of interest. This mirrors physical
environments, where cities scale by adding neighborhoods rather than forcing all activity
through a single thoroughfare.

This spatialization fundamentally alters adversarial scaling dynamics. Engagement--driven
platforms monetize attention by selling advertising inventory, frequently to actors whose
business models depend on deception, including known scam networks. These platforms
therefore benefit financially from maximizing exposure, regardless of semantic quality.
In contrast, a spatial semantic medium offers no global amplification mechanism to exploit.
Content remains localized unless structurally integrated into broader semantic regions.

\subsection{Onboarding Through Programming as Literacy}

A frequent critique of structured systems is their perceived difficulty for newcomers.
However, the platform under consideration explicitly incorporates programming education as
a first--class function through the SpherePOP calculus. Users are not expected to master
formal systems upfront; rather, they learn by interacting with semantic objects whose
behavior is governed by simple, composable rules.

Programming in this context is not an auxiliary skill but a literacy: the ability to
construct, transform, and navigate meaning. By embedding programming instruction directly
into social interaction, the system converts what is traditionally a barrier into a
continuous learning process.

This contrasts sharply with existing platforms, where users are incentivized to spend
substantial time curating follow lists, monitoring engagement metrics, and pursuing
arbitrary follower goals. Such activity consumes attention without producing durable
knowledge or skill. In the coherence--first system, time invested in learning the substrate
directly increases one’s expressive and navigational capacity.

\subsection{Persistence, Memory, and Reaccessibility}

Semantic persistence introduces a distinct limitation: content cannot be trivially
forgotten. However, this persistence enables a powerful alternative to feed consumption.
Knowledge is stored in stable locations within the semantic latent space and can be
reaccessed by returning to the same region.

This property transforms memory from a chronological archive into a spatial one. Users do
not scroll backward through time but revisit conceptual neighborhoods. Over time, these
neighborhoods accrete structure, becoming richer rather than buried.

While this challenges habitual consumption patterns shaped by infinite scroll interfaces,
it aligns more closely with how expertise and understanding develop. The cost is reduced
novelty throughput; the benefit is cumulative intelligibility.

\subsection{Adversarial Behavior in a Navigable Space}

Adversarial behavior in a spatial semantic system differs qualitatively from that in
broadcast platforms. Without feeds or engagement metrics, attackers cannot leverage
visibility cascades. Instead, malicious content must occupy space, persist, and withstand
entropy constraints.

Attempts to game entropy budgets or flood regions with low--quality objects are constrained
by both computational and navigational friction. Users encounter such content only by
entering the relevant regions, and repeated incoherent transformations accumulate entropy
that eventually blocks further evolution.

Crucially, the absence of an advertising--driven revenue model removes the primary
financial incentive for large--scale spam. Traditional platforms profit from keeping users
engaged, even if that engagement consists of navigating scams, outrage, or empty status
games. A spatial semantic platform derives no benefit from time--wasting behavior oriented
around following, metric optimization, or performative presence.

\subsection{Expression, Revision, and Spatial Accountability}

Freedom of expression in a spatial semantic medium is preserved through transformation
rather than erasure. Users may revise, retract, or contextualize prior contributions, but
these actions are recorded as movements or transformations within semantic space.

This introduces a form of spatial accountability. Past statements remain accessible at
their original locations, while newer interpretations may occupy adjacent or higher--order
regions. Growth is represented geometrically rather than narratively.

While this model may feel restrictive to users accustomed to disappearing content, it
reflects norms already present in scholarly and technical domains, where retracted or
superseded work remains part of the record.

\subsection{Migration and Hybrid Use}

Migration from existing platforms remains nontrivial, but the spatial model enables hybrid
strategies. External content can be imported as semantic objects anchored to regions of
space that reflect their topical and structural properties. Over time, these imports may
be extended, formalized, or abandoned.

Notably, the system does not attempt to import engagement artifacts such as likes, shares,
or follower counts. These signals have no spatial or semantic meaning and are therefore
discarded. Users migrating into the system are freed from the obligation to maintain
continuous visibility or pursue follower accumulation as a proxy for participation.

\subsection{Summary of Tradeoffs}

The coherence--first, spatial semantic platform prioritizes cumulative understanding,
programming literacy, and persistent memory over immediacy and attention capture. It is
therefore incompatible with advertising--driven social media models and the behavioral
patterns they encourage.

Its limitations are real: increased cognitive demand, slower initial growth, and a
requirement that users engage actively with structure. Its advantages are equally real:
durable knowledge, teachable computation, and a social medium in which time spent produces
skill and understanding rather than exhaustion.

\section{Formalization and Quantitative Analysis}

This section formalizes the central claims of the coherence--first spatial semantic
architecture and analyzes its properties using tools from information theory, dynamical
systems, and computational complexity. We compare feed--based dissemination with semantic
neighborhood navigation, model trust accumulation under persistence, and establish
convergence guarantees under entropy--bounded evolution.

\subsection{Preliminaries}

Let $\Sigma$ denote the space of all semantic objects, and let $X$ denote the semantic
latent manifold in which objects are embedded. Each semantic object $\sigma \in \Sigma$
is associated with an embedding $\phi(\sigma) \in X$, a provenance graph $P(\sigma)$, and
a semantic entropy value $E(\sigma)$.

Users interact with the system by navigating $X$, applying transformations to semantic
objects, and composing programs in a structured transformation language.

\subsection{Feed--Based Dissemination versus Semantic Neighborhoods}

We first contrast feed--based dissemination with spatial semantic navigation.

\begin{definition}[Feed Model]
A feed--based system presents users with a sequence
\[
F_u(t) = \langle c_1, c_2, \dots, c_n \rangle
\]
ordered by a relevance function $R(c, u, t)$ optimized for engagement metrics such as
clicks, likes, or dwell time.
\end{definition}

\begin{definition}[Semantic Neighborhood Model]
A semantic neighborhood system presents users with content
\[
N_u(x, r) = \{ \sigma \in \Sigma \mid d(\phi(\sigma), x) \le r \}
\]
where $x \in X$ is the user’s current location in semantic space and $r$ is a bounded
radius.
\end{definition}

\begin{proposition}[Entropy Accumulation in Feeds]
Feed--based dissemination induces unbounded semantic entropy accumulation over time.
\end{proposition}

\begin{proof}[Sketch]
In feed systems, content relevance decays with time and is replaced by new content
optimized for engagement rather than coherence. There is no constraint enforcing semantic
consistency across items in $F_u(t)$. As $t \to \infty$, the expected semantic divergence
between consecutive items grows without bound, yielding unbounded entropy.
\end{proof}

\begin{proposition}[Entropy Localization in Semantic Neighborhoods]
In a semantic neighborhood system, entropy accumulation is locally bounded.
\end{proposition}

\begin{proof}[Sketch]
Semantic neighborhoods restrict interaction to bounded regions of $X$. Transformations
applied within a neighborhood must satisfy entropy budgets. Therefore, entropy growth is
constrained locally and does not propagate globally.
\end{proof}

\subsection{Information--Theoretic Comparison}

Let $H(F_u)$ denote the entropy of a user’s feed and $H(N_u)$ the entropy of a semantic
neighborhood.

\begin{theorem}[Asymptotic Entropy Divergence]
For feed--based systems,
\[
\lim_{t \to \infty} H(F_u(t)) = \infty,
\]
whereas for semantic neighborhoods,
\[
\sup_t H(N_u(x,t)) \le H_{\max}
\]
for some finite bound $H_{\max}$.
\end{theorem}

\begin{proof}[Sketch]
Feed entropy increases as content turnover introduces uncorrelated items optimized for
engagement. Semantic neighborhoods impose correlation through spatial proximity and
entropy--bounded transformations, yielding a finite upper bound.
\end{proof}

This result formalizes the intuition that feeds maximize novelty while neighborhoods
preserve structure.

\subsection{Trust Accumulation: Persistence versus Popularity}

We now model trust accumulation.

\begin{definition}[Popularity--Based Trust]
In engagement systems, trust is approximated by a function
\[
T_{\text{pop}}(u,t) = f(L_u(t), S_u(t), F_u(t))
\]
where $L$ denotes likes, $S$ shares, and $F$ follower count.
\end{definition}

\begin{definition}[Persistence--Based Trust]
In a coherence--first system, trust is defined as
\[
T_{\text{pers}}(u,t) = g(|P_u(t)|, \Delta E_u(t), C_u(t))
\]
where $|P_u|$ is the size of the user’s provenance graph, $\Delta E_u$ measures entropy
stability over time, and $C_u$ denotes successful reconciliations.
\end{definition}

\begin{proposition}[Instability of Popularity Trust]
Popularity--based trust is unstable under adversarial optimization.
\end{proposition}

\begin{proof}[Sketch]
Engagement metrics can be manipulated through coordinated activity, bot networks, or
sensational content. Small perturbations in visibility can cause large changes in
$T_{\text{pop}}$, rendering it unreliable.
\end{proof}

\begin{theorem}[Monotonicity of Persistence Trust]
Persistence--based trust is monotonic under valid transformations.
\end{theorem}

\begin{proof}
Valid transformations append to provenance and are entropy--bounded. Therefore, $|P_u|$
and reconciliation counts increase monotonically, while entropy stability penalizes
destabilizing behavior.
\end{proof}

\subsection{Convergence and Stability under Entropy Bounds}

We model semantic evolution as a dynamical system.

\begin{definition}[Semantic Evolution Operator]
Let $\mathcal{T}$ denote the set of admissible transformations. Define
\[
\mathcal{E}_{t+1} = \mathcal{T}(\mathcal{E}_t)
\]
subject to entropy constraints.
\end{definition}

\begin{theorem}[Convergence of Semantic Objects]
Under fixed entropy budgets, semantic object evolution converges to stable attractors in
$\Sigma$.
\end{theorem}

\begin{proof}[Sketch]
Entropy budgets impose a Lyapunov--like constraint on semantic evolution. Since entropy
cannot increase without bound and transformations are finite, sequences converge to fixed
points or bounded cycles representing stable interpretations.
\end{proof}

This establishes long--term stability absent from feed--based systems.

\subsection{Programming Literacy and Learning Efficiency}

Let $L_u(t)$ denote a user’s programming literacy as measured by successful semantic
transformations composed in the transformation language.

\begin{proposition}[Positive Learning Gradient]
In a system where social interaction is mediated through programmable transformations,
expected programming literacy increases monotonically with use.
\end{proposition}

\begin{proof}[Sketch]
Each interaction requires the application or observation of typed transformations. Unlike
feed consumption, which yields no cumulative skill, repeated interaction increases the
user’s effective action space.
\end{proof}

This contrasts with engagement systems, where time spent pursuing follower goals yields no
durable competence.

\subsection{Computational Complexity}

We conclude with complexity analysis.

Let $n$ be the number of semantic objects, $k$ the maximum degree of provenance graphs,
and $m$ the number of transformations.

\begin{proposition}[Transformation Complexity]
Applying a transformation has time complexity $O(k)$ and space complexity $O(1)$.
\end{proposition}

\begin{proposition}[Neighborhood Query Complexity]
Semantic neighborhood queries execute in $O(\log n + |N_u|)$ using standard spatial
indexing.
\end{proposition}

\begin{theorem}[Global Scalability]
Total system complexity grows sublinearly with user count under bounded neighborhood
interaction.
\end{theorem}

\begin{proof}[Sketch]
Users interact locally in semantic space. No global feed requires recomputation or ranking
across all objects. Therefore, system load scales with local density rather than total
population.
\end{proof}

This completes the formal and quantitative analysis of the coherence--first semantic
architecture.

\section{Artificial Intelligence as a Typed Semantic Operator}

In the coherence--first semantic architecture, artificial intelligence is neither a
primary agent nor a curator of attention. Instead, AI functions as a constrained operator
within the semantic substrate, performing well--defined transformations under explicit
structural and entropy constraints. This section formalizes the role of AI and contrasts
it with the recommender--centric models that dominate contemporary social media.

\subsection{Rejection of AI as a Relevance Oracle}

Most social media platforms deploy AI primarily as a relevance oracle: a system that
predicts which content will maximize user engagement. This framing aligns AI incentives
with time spent, emotional arousal, and interaction volume. As a result, recommender
systems become amplifiers of whatever content best exploits human cognitive biases,
regardless of semantic quality.

The coherence--first architecture explicitly rejects this role. AI systems are not tasked
with ranking users, optimizing feeds, or predicting attention. There is no global feed to
optimize, and no engagement metric to maximize. Consequently, entire classes of alignment
failures associated with recommender systems are rendered inapplicable.

\subsection{Typed Semantic Operations}

AI operates exclusively through typed semantic transformations.

\begin{definition}[Typed Semantic Operator]
A typed semantic operator is a transformation
\[
T : (\Sigma, \mathcal{M}_{\text{in}}) \rightarrow (\Sigma, \mathcal{M}_{\text{out}})
\]
that maps semantic objects with input modality requirements $\mathcal{M}_{\text{in}}$ to
objects with output modality requirements $\mathcal{M}_{\text{out}}$, subject to fixed
entropy budgets.
\end{definition}

Examples include summarization, translation, modality completion, formalization,
compression, and reconciliation assistance. Each operator declares the modalities it
consumes and produces, enabling static validation prior to execution.

\subsection{Entropy Budgets and Validation}

Every AI operator is associated with an entropy budget $\varepsilon_T$. After application,
the resulting semantic object is evaluated for coherence.

\begin{proposition}[Post--Transformation Validation]
An AI--generated transformation is accepted if and only if the resulting semantic object
is well--formed and satisfies the entropy bound
\[
E(T(\sigma)) \le E(\sigma) + \varepsilon_T.
\]
\end{proposition}

If validation fails, the transformation is rejected or returned for human mediation. AI
therefore cannot unilaterally inject ambiguity, contradiction, or scope inflation into the
system.

\subsection{Failure Modes and Recovery}

AI transformations may fail due to insufficient input structure, modality mismatch, or
exceeding entropy budgets. Rather than silently degrading output quality, such failures are
explicitly surfaced.

Failed transformations generate diagnostic artifacts indicating missing modalities,
uncertain mappings, or unresolved inconsistencies. These artifacts become semantic objects
in their own right, enabling users to inspect and address the failure.

This contrasts with engagement systems, where AI failures are often invisible and manifest
only indirectly through degraded discourse quality.

\subsection{Training Data and Bias Mitigation}

Because AI operators perform constrained transformations rather than open--ended
generation, training data requirements differ substantially from those of large language
models deployed as conversational agents. Training focuses on mapping between modalities
and preserving semantic invariants rather than imitating human conversational behavior.

Bias mitigation is addressed structurally. Since AI output must satisfy explicit semantic
constraints and provenance requirements, biased or misleading transformations are more
readily detectable. Moreover, the absence of engagement optimization removes incentives to
exploit polarizing or sensational patterns present in training data.

\subsection{Human--AI Collaboration}

AI in a coherence--first system functions as a collaborator rather than a substitute for
human judgment. It assists users in navigating semantic space, identifying inconsistencies,
and performing routine transformations, but it does not author meaning independently.

This role aligns AI incentives with clarity and stability rather than persuasion or
attention capture. By constraining AI to typed operations with explicit budgets, the system
ensures that human agency and accountability remain central.

\subsection{Summary}

By redefining AI as a typed semantic operator rather than a relevance oracle, the
coherence--first architecture avoids the alignment failures endemic to engagement--driven
systems. AI becomes a tool for maintaining semantic integrity and supporting learning,
including programming literacy, rather than a mechanism for amplifying noise or optimizing
visibility.

\section{Governance, Invariants, and Forking Mechanics}

Governance in feed--based social media platforms is typically implemented as an external
layer: policies, moderators, enforcement actions, and appeals. These mechanisms operate
reactively, intervening after semantic damage has occurred and often at scales that render
consistent enforcement impractical. In contrast, the coherence--first semantic
architecture embeds governance into the formal structure of the system itself. This
section specifies the invariants that define admissible behavior, the procedural handling
of conflict and divergence, and the mechanics of forking as a first--class operation.

\subsection{Governance by Structural Invariants}

At the core of the system is a set of invariants that all valid semantic states must
satisfy. These invariants are enforced mechanically rather than socially.

\begin{definition}[Semantic Invariant]
A semantic invariant is a property $I(\sigma)$ of a semantic object $\sigma$ that must
hold for $\sigma$ to be considered valid within the system.
\end{definition}

Examples include modality completeness, provenance monotonicity, and bounded semantic
entropy. Violations of invariants do not trigger punitive action; instead, they prevent
the object from participating in further transformations until resolved.

\begin{proposition}[Invariant--Based Governance]
All admissible system states satisfy the set of semantic invariants.
\end{proposition}

\begin{proof}
Invariant checks are applied at the time of transformation. Since invalid states are
rejected by construction, no reachable system state violates the invariants.
\end{proof}

This approach shifts governance from discretionary enforcement to constraint satisfaction.

\subsection{Semantic Reconciliation}

Disagreement is an expected feature of discourse. The coherence--first system does not
attempt to eliminate disagreement but to represent it explicitly.

\begin{definition}[Semantic Reconciliation]
Semantic reconciliation is a transformation that integrates divergent branches of a
semantic object by introducing additional structure that resolves or contextualizes their
differences.
\end{definition}

Reconciliation may involve narrowing scope, introducing conditionality, or adding
supporting modalities. Crucially, reconciliation is optional rather than mandatory; the
system permits persistent divergence when resolution is not possible or desirable.

\begin{proposition}[Non--Coercive Resolution]
No semantic object is required to reconcile divergent branches.
\end{proposition}

\begin{proof}
Branching is a valid terminal state. The absence of a reconciliation transformation does
not violate any invariant.
\end{proof}

This ensures that governance does not collapse into enforced consensus.

\subsection{Forking as a First--Class Operation}

Forking is treated as a constructive operation rather than a failure mode.

\begin{definition}[Semantic Fork]
A semantic fork is the creation of two or more descendant semantic objects that diverge
from a common ancestor due to incompatible transformations.
\end{definition}

Forks preserve shared provenance up to the point of divergence and then evolve
independently. Unlike moderation--based deletion or suppression, forking externalizes
conflict rather than obscuring it.

\subsection{Tracking Divergence}

Divergence is tracked explicitly in the provenance graph. Users navigating semantic space
can observe where and why forks occurred, inspect the transformations that produced them,
and choose which branches to engage with.

This transparency contrasts with feed--based systems, where divergence is often masked by
algorithmic filtering or visibility dynamics.

\subsection{Coordinated Attacks and Structural Resistance}

Coordinated attacks, such as brigading or disinformation campaigns, exploit amplification
mechanisms in engagement--driven systems. The coherence--first architecture removes these
mechanisms entirely.

Without feeds, likes, or shares, attackers cannot leverage visibility cascades. Malicious
content must persist, occupy semantic space, and withstand invariant checks. Coordinated
low--quality transformations accumulate entropy and exhaust available budgets, limiting
their impact.

\subsection{Emergent Norms}

While formal invariants define the boundaries of admissible behavior, norms emerge through
use. Because semantic objects persist and interactions are traceable, norms stabilize over
time rather than being enforced episodically.

This separation of formal constraints from emergent norms allows the system to accommodate
diverse communities without imposing uniform standards beyond those required for semantic
coherence.

\subsection{Summary}

Governance in a coherence--first semantic substrate is achieved through invariants,
transparent reconciliation, and explicit forking rather than moderation and suppression.
By embedding governance into the structure of meaning itself, the system avoids many of
the failure modes associated with discretionary enforcement and algorithmic opacity.

\section{Implementation Considerations}

While the coherence--first semantic architecture is motivated by theoretical and
information--theoretic considerations, its viability depends on concrete implementation
choices. This section outlines practical considerations related to storage, identity,
networking, and performance. The goal is not to prescribe a single implementation, but to
demonstrate that the proposed architecture is computationally feasible with existing
technology.

\subsection{Storage of Semantic Objects and Provenance}

Semantic objects are persistent, immutable entities whose evolution is recorded through
append--only provenance graphs. This naturally suggests a storage model based on
content--addressed data structures. Each semantic object state can be identified by a
cryptographic hash of its content, modality assignments, and provenance references.

Provenance graphs grow monotonically but typically exhibit low branching factors in
practice. Since transformations are localized and entropy--bounded, most objects evolve
through short, interpretable chains rather than deep or highly branching structures. This
enables efficient storage through structural sharing, similar to persistent data
structures used in functional programming languages.

Garbage collection is unnecessary in the traditional sense, as objects are never mutated
or deleted. Instead, inactive or superseded objects can be indexed separately or cached at
lower priority without violating persistence guarantees.

\subsection{Indexing and Semantic Space Navigation}

Efficient navigation of semantic neighborhoods requires spatial indexing over the latent
manifold $X$. Standard techniques such as approximate nearest--neighbor search, metric
trees, or locality--sensitive hashing can be employed to support neighborhood queries.

Because users interact locally within semantic space, query complexity depends primarily
on local density rather than global object count. This property enables horizontal scaling
by partitioning semantic space across nodes while preserving navigability.

Importantly, indexing operates over semantic embeddings rather than engagement metrics.
This eliminates the need for continuous global re-ranking and reduces system--wide
recomputation.

\subsection{Identity Anchoring and Authentication}

Identity in the coherence--first system is defined by persistent identifiers rather than
profiles or accounts. Cryptographic primitives can be used to anchor identity securely.
Each identity may be associated with a public key, and transformations can be signed to
establish authorship without revealing unnecessary personal information.

This model supports pseudonymity while preserving accountability. Because identity is
anchored structurally rather than socially, reputation laundering through account
recreation or metric manipulation is ineffective.

\subsection{Distributed Deployment and Federation}

The architecture is compatible with distributed and federated deployment. Semantic space
can be partitioned across nodes or instances, with each node responsible for maintaining a
region of the latent manifold and the objects embedded within it.

Federation does not require global agreement on ranking or moderation policies. Nodes need
only agree on core invariants governing semantic object validity and transformation. This
reduces coordination overhead and allows heterogeneous communities to coexist without
fragmenting the substrate.

Cross--node interaction is mediated through provenance references and embedding mappings,
ensuring that semantic integrity is preserved even when objects traverse administrative
boundaries.

\subsection{Performance Characteristics}

Performance in engagement--driven platforms is dominated by feed ranking and recommendation
pipelines, which require continuous large--scale computation. In contrast, the
coherence--first architecture shifts computation toward transformation validation and
spatial queries.

Transformation validation is localized and bounded in complexity, as entropy budgets and
modality checks constrain the scope of evaluation. Spatial queries benefit from standard
indexing techniques and do not require system--wide recomputation.

As a result, system load scales with active semantic regions rather than total user count
or content volume. This property supports long--term growth without the runaway
computational costs associated with global feeds.

\subsection{Programming Interface and SpherePOP Integration}

The integration of a programmable transformation language, such as SpherePOP, introduces
additional implementation considerations. Programs operate over semantic objects and must
respect modality and entropy constraints. Static analysis can be employed to verify that
programs are well--typed and cannot violate invariants.

Because programming constructs are embedded directly into social interaction, tooling for
debugging, visualization, and explanation is essential. However, the simplicity of the
underlying calculus enables incremental learning and compositional reuse.

By treating programming as a medium of interaction rather than a specialized activity,
the system aligns implementation complexity with user skill development rather than
hiding it behind opaque automation.

\subsection{Summary}

The coherence--first semantic architecture can be implemented using existing storage,
indexing, cryptographic, and distributed systems techniques. Its performance profile
differs substantially from that of engagement--driven platforms, emphasizing locality,
persistence, and validation over global optimization. These properties make it suitable
for scalable deployment without reproducing the incentive structures that undermine
contemporary social media.

\section{Conclusion and Research Agenda}

This paper has argued that the dominant failures of contemporary social media platforms
are not incidental defects, but necessary consequences of their underlying architecture.
Feed--based dissemination, engagement optimization, and performative identity together
produce high--entropy dynamics that degrade meaning, incentivize time--wasting behavior,
and erode trust at scale. Attempts to repair these systems through moderation, transparency,
or algorithmic adjustment leave the core substrate unchanged and therefore fail to address
the root causes.

In response, we have outlined a coherence--first alternative grounded in persistent
semantic objects, entropy--bounded evolution, provenance--anchored identity, and spatial
navigation of meaning. Rather than treating communication as a stream of posts optimized
for visibility, the proposed architecture treats meaning as a structured entity that
evolves through constrained transformations. Trust emerges from persistence and
consistency over time rather than popularity or attention metrics.

We formalized the key components of this architecture, including semantic objects,
transformation operators, entropy budgets, and provenance graphs. We demonstrated that
feed--based systems exhibit unbounded entropy accumulation, while semantic neighborhoods
localize and constrain disorder. We showed that trust models based on persistence are
structurally more stable than those based on popularity, and that entropy--bounded
evolution yields convergence properties absent from engagement--driven platforms. We
further analyzed computational complexity and demonstrated that the system scales through
locality rather than global optimization.

Importantly, the architecture reframes the role of artificial intelligence. By constraining
AI to typed semantic operations under explicit budgets, the system avoids the alignment
failures associated with recommender systems and attention optimization. AI becomes a tool
for maintaining coherence, assisting learning, and supporting programming literacy rather
than amplifying noise or persuasion.

The integration of a programmable semantic calculus transforms social interaction into a
learning process. Time spent within the system produces durable skills and knowledge rather
than ephemeral engagement. Knowledge persists spatially and can be reaccessed by returning
to the same region of semantic space, enabling cumulative understanding rather than
chronological decay.

\subsection{Open Problems and Future Work}

Several open problems remain. First, while entropy--bounded evolution provides strong
stability guarantees, further work is needed to refine entropy metrics that balance
flexibility with constraint across diverse domains. Second, the design of user interfaces
for high--dimensional semantic navigation remains an active area of research, particularly
with respect to accessibility and cognitive load. Third, empirical studies are required to
evaluate learning outcomes, trust formation, and discourse quality in deployed systems.

Additional work is needed to explore incentive models compatible with coherence--first
design. Advertising--driven revenue is fundamentally misaligned with persistence and
semantic integrity. Alternative models, such as patronage, institutional support, or
public infrastructure funding, warrant systematic investigation.

Finally, the relationship between coherence--first platforms and existing digital
ecosystems remains to be fully articulated. Hybrid deployments, migration strategies, and
interoperability with legacy systems present both technical and social challenges.

\subsection{Broader Implications}

Beyond social media, the coherence--first architecture has implications for knowledge
management, education, scientific collaboration, and governance. Any domain in which
meaning must persist, evolve, and remain accountable over time stands to benefit from a
substrate that prioritizes structure over visibility.

By rethinking the primitives of digital interaction, this work contributes to a growing
body of research that treats platform design as a problem of semantics and systems
architecture rather than content moderation alone. It suggests that healthier digital
public spheres require not better incentives layered atop existing systems, but new
substrates in which certain failure modes are structurally impossible.

\subsection{Closing Remarks}

The central claim of this paper is that coherence is not an emergent property of
engagement--optimized systems, but a design choice that must be encoded at the substrate
level. A coherence--first semantic architecture offers a viable alternative to social
media as currently constituted, one in which time spent produces understanding rather than
exhaustion, and where trust is earned through persistence rather than performance.

We invite further research, experimentation, and critique of this approach as part of a
broader effort to reimagine the foundations of digital communication.

\newpage

\section*{Appendices} 

\appendix
\section{Formal Semantic Object Model}

\subsection{Primitive Sets}

Let:
\[
\mathcal{I} \text{ be a countable set of immutable identifiers}
\]
\[
\mathcal{M} \text{ be a finite set of modality labels}
\]
\[
\mathcal{C} \text{ be a space of modality-indexed content realizations}
\]
\[
\mathcal{P} \text{ be the space of finite directed acyclic graphs}
\]
\[
\mathbb{R}_{\ge 0} \text{ be the non-negative reals}
\]

\subsection{Semantic Object Definition}

\begin{definition}[Semantic Object]
A semantic object is a tuple
\[
\sigma = (i, \mathcal{M}_\sigma, C_\sigma, P_\sigma, E_\sigma)
\]
where:
\begin{align*}
i &\in \mathcal{I} \quad \text{(immutable identity)} \\
\mathcal{M}_\sigma &\subseteq \mathcal{M} \quad \text{(required modalities)} \\
C_\sigma &: \mathcal{M}_\sigma \rightharpoonup \mathcal{C} \quad \text{(partial content map)} \\
P_\sigma &\in \mathcal{P} \quad \text{(provenance DAG)} \\
E_\sigma &\in \mathbb{R}_{\ge 0} \quad \text{(semantic entropy)}
\end{align*}
\end{definition}

\subsection{Well-Formedness}

\begin{definition}[Modal Completeness]
A semantic object $\sigma$ is modally complete iff:
\[
\forall m \in \mathcal{M}_\sigma,\; C_\sigma(m) \downarrow
\]
where $\downarrow$ denotes definedness.
\end{definition}

\begin{definition}[Provenance Validity]
A provenance graph $P_\sigma = (V,E)$ is valid iff:
\[
\text{$P_\sigma$ is finite, directed, and acyclic}
\]
\end{definition}

\begin{definition}[Well-Formed Semantic Object]
A semantic object $\sigma$ is well-formed iff:
\[
\text{ModalCompleteness}(\sigma) \;\wedge\; \text{ProvenanceValidity}(P_\sigma)
\]
\end{definition}

\subsection{Semantic State Space}

\begin{definition}[Semantic State Space]
Define the semantic state space as:
\[
\Sigma = \{ \sigma \mid \sigma \text{ is well-formed} \}
\]
\end{definition}

\subsection{Identity Invariance}

\begin{definition}[Identity Preservation]
For any transformation $T : \Sigma \to \Sigma$,
\[
T(\sigma) = \sigma' \implies i_{\sigma'} = i_\sigma
\]
\end{definition}

\begin{theorem}[Identity Invariance]
Identity is invariant under all admissible transformations.
\end{theorem}

\begin{proof}
By definition, transformations operate on content, provenance, and entropy only. The
identity component is immutable and not in the transformation domain.
\end{proof}

---

\subsection{Content Extensionality}

\begin{definition}[Extensional Equality]
Two semantic objects $\sigma_1, \sigma_2$ are extensionally equal iff:
\[
i_1 = i_2 \;\wedge\; C_{\sigma_1} = C_{\sigma_2} \;\wedge\; P_{\sigma_1} = P_{\sigma_2}
\]
\end{definition}

\begin{proposition}[No Implicit Overwrite]
If $\sigma_1 \neq \sigma_2$ extensionally, then $\sigma_1$ cannot overwrite $\sigma_2$.
\end{proposition}

\subsection{Structural Consequences}

\begin{theorem}[No Atomic Posts]
There exists no semantic object with empty provenance and mutable identity.
\end{theorem}

\begin{proof}
Provenance monotonicity and identity immutability prohibit atomic overwrite semantics.
\end{proof}

\begin{theorem}[No Context-Free Meaning]
There exists no $\sigma \in \Sigma$ such that $C_\sigma$ is defined independently of
$P_\sigma$.
\end{theorem}

\subsection{Remarks}

All social and computational interaction in the system is reducible to operations over
$\Sigma$ subject to the constraints defined above.

\section{Semantic Entropy Measure and Bounds}

\subsection{Preliminaries}

Let $\Sigma$ denote the space of well-formed semantic objects defined in Appendix~A.
For $\sigma \in \Sigma$, recall:
\[
\sigma = (i, \mathcal{M}_\sigma, C_\sigma, P_\sigma, E_\sigma).
\]

Let $\Omega$ be a measurable space of semantic interpretations, and let
\[
\mu_\sigma : \Omega \to [0,1]
\]
be a probability measure induced by the content and provenance of $\sigma$.

\subsection{Semantic Entropy Definition}

\begin{definition}[Semantic Entropy Functional]
Define semantic entropy as:
\[
E(\sigma) := H(\mu_\sigma) = - \int_\Omega \mu_\sigma(\omega)\log \mu_\sigma(\omega)\, d\omega
\]
where $H$ denotes Shannon entropy.
\end{definition}

\begin{definition}[Conditional Semantic Entropy]
Given provenance $P_\sigma$, define:
\[
E(\sigma \mid P_\sigma) := H(\mu_\sigma \mid P_\sigma).
\]
\end{definition}

This quantity measures ambiguity remaining after accounting for derivational context.

\subsection{Entropy Decomposition}

\begin{lemma}[Entropy Decomposition]
For any $\sigma \in \Sigma$,
\[
E(\sigma) = E(C_\sigma) + E(P_\sigma) - I(C_\sigma; P_\sigma)
\]
where $I$ is mutual information.
\end{lemma}

\begin{proof}
Direct from Shannon entropy identities.
\end{proof}

\subsection{Entropy-Bounded Transformations}

\begin{definition}[Transformation Operator]
A semantic transformation is a function:
\[
T : \Sigma \to \Sigma
\]
such that identity and provenance monotonicity are preserved.
\end{definition}

\begin{definition}[Entropy Budget]
Each transformation $T$ is associated with a fixed constant $\varepsilon_T \ge 0$.
\end{definition}

\begin{definition}[Entropy-Bounded Transformation]
$T$ is entropy-bounded iff:
\[
E(T(\sigma)) \le E(\sigma) + \varepsilon_T.
\]
\end{definition}

\subsection{Additivity and Subadditivity}

\begin{lemma}[Entropy Subadditivity]
For a sequence of transformations $(T_k)$ applied to $\sigma_0$,
\[
E(\sigma_n) \le E(\sigma_0) + \sum_{k=1}^n \varepsilon_{T_k}.
\]
\end{lemma}

\begin{proof}
By induction on $n$ using the entropy-bounded condition.
\end{proof}

\subsection{Lyapunov Stability}

\begin{definition}[Semantic Lyapunov Function]
Define $V(\sigma) := E(\sigma)$.
\end{definition}

\begin{theorem}[Semantic Stability]
If $\sum_k \varepsilon_{T_k} < \infty$, then $(\sigma_k)$ converges in entropy.
\end{theorem}

\begin{proof}
$V(\sigma_k)$ is monotonically bounded above and hence convergent.
\end{proof}

\subsection{Entropy Laundering Impossibility}

\begin{definition}[Entropy Laundering Attempt]
An entropy laundering attempt is a finite sequence of transformations $(T_k)$ such that:
\[
E(\sigma_n) < E(\sigma_0)
\]
while increasing semantic ambiguity.
\end{definition}

\begin{theorem}[No Entropy Laundering]
Entropy laundering is impossible under entropy-bounded transformations.
\end{theorem}

\begin{proof}
Entropy reduction requires negative $\varepsilon_T$, which is disallowed. Any reduction in
$E(\sigma)$ must arise from added structure or information, not reordering or deletion.
\end{proof}

---

\subsection{Branching and Entropy}

\begin{lemma}[Branch Entropy Conservation]
Let $\sigma \to \{\sigma_1, \sigma_2\}$ be a fork. Then:
\[
E(\sigma_1) + E(\sigma_2) \ge E(\sigma).
\]
\end{lemma}

\begin{proof}
Branching introduces independent uncertainty components; entropy is not destroyed.
\end{proof}

\subsection{Entropy Caps}

\begin{definition}[Global Entropy Cap]
A system-wide constant $E_{\max}$ such that:
\[
E(\sigma) \le E_{\max} \quad \forall \sigma \in \Sigma.
\]
\end{definition}

\begin{theorem}[Structural Drift Prevention]
Unbounded semantic drift is structurally impossible if $E_{\max} < \infty$.
\end{theorem}

\begin{proof}
Immediate from boundedness of $V(\sigma)$.
\end{proof}

\subsection{Consequences}

\begin{corollary}[No Viral Collapse]
No semantic object can accumulate arbitrarily high ambiguity via repeated exposure.
\end{corollary}

\begin{corollary}[No Engagement-Driven Explosion]
Visibility amplification without semantic structure cannot increase entropy budget.
\end{corollary}

\section{Provenance Graph Structure}

\subsection{Graph-Theoretic Preliminaries}

Let $\mathcal{G}$ denote the class of finite directed graphs.

\begin{definition}[Directed Acyclic Graph]
A directed graph $G = (V,E)$ is a DAG iff:
\[
\forall v \in V,\; \nexists \text{ directed cycle through } v.
\]
\end{definition}

Let $\prec$ denote the strict partial order induced by reachability in a DAG.

\subsection{Provenance Graph Definition}

\begin{definition}[Provenance Graph]
For a semantic object $\sigma$, its provenance graph is
\[
P_\sigma = (V_\sigma, E_\sigma)
\]
where:
\begin{align*}
V_\sigma &\subset \Sigma \quad \text{(semantic states)} \\
E_\sigma &\subset V_\sigma \times V_\sigma \quad \text{(typed transformations)}
\end{align*}
and $P_\sigma$ is a DAG.
\end{definition}

\subsection{Causal Ordering}

\begin{definition}[Causal Precedence]
For $\sigma_i, \sigma_j \in V_\sigma$,
\[
\sigma_i \prec \sigma_j \iff \exists \text{ directed path } \sigma_i \to \sigma_j.
\]
\end{definition}

\begin{lemma}[Irreflexivity]
\[
\neg(\sigma \prec \sigma)
\]
\end{lemma}

\begin{lemma}[Transitivity]
\[
(\sigma_i \prec \sigma_j) \wedge (\sigma_j \prec \sigma_k) \Rightarrow (\sigma_i \prec \sigma_k).
\]
\end{lemma}

Thus $\prec$ defines a strict partial order.

\subsection{Monotonicity Invariant}

\begin{definition}[Provenance Monotonicity]
A provenance graph satisfies monotonicity iff:
\[
\forall (\sigma_i \to \sigma_j) \in E_\sigma,\quad
V_{\sigma_i} \subset V_{\sigma_j}.
\]
\end{definition}

\begin{theorem}[Monotonic Growth]
Provenance graphs grow monotonically under all admissible transformations.
\end{theorem}

\begin{proof}
Transformations append nodes and edges without deletion by definition.
\end{proof}

\subsection{Forking}

\begin{definition}[Fork]
A fork occurs when a node $\sigma$ has two outgoing edges:
\[
\sigma \to \sigma_1,\quad \sigma \to \sigma_2
\]
such that $\sigma_1 \not\prec \sigma_2$ and $\sigma_2 \not\prec \sigma_1$.
\end{definition}

\begin{lemma}[Fork Non-Equivalence]
Forked descendants are not extensionally equal:
\[
\sigma_1 \neq \sigma_2.
\]
\end{lemma}

\subsection{Merge (Reconciliation)}

\begin{definition}[Merge Operator]
A merge is a transformation
\[
M : (\sigma_1, \sigma_2) \mapsto \sigma_m
\]
such that:
\[
\sigma_1 \prec \sigma_m \quad \wedge \quad \sigma_2 \prec \sigma_m.
\]
\end{definition}

\begin{definition}[Merge Validity]
A merge is valid iff:
\[
E(\sigma_m) \le \max(E(\sigma_1), E(\sigma_2)) + \varepsilon_M.
\]
\end{definition}

\subsection{History Preservation}

\begin{theorem}[No History Erasure]
There exists no admissible transformation $T$ such that:
\[
P_{T(\sigma)} \subsetneq P_\sigma.
\]
\end{theorem}

\begin{proof}
Such a transformation would violate monotonicity and acyclicity.
\end{proof}

\subsection{Non-Linear Time}

\begin{definition}[Semantic Time]
Semantic time is defined by the partial order $(V_\sigma, \prec)$, not by linear indices.
\end{definition}

\begin{proposition}[No Total Order]
There exists no total order compatible with $\prec$ for all provenance graphs.
\end{proposition}

\begin{proof}
Forks induce incomparable nodes.
\end{proof}

\subsection{Provenance Isomorphism}

\begin{definition}[Provenance Isomorphism]
Two provenance graphs $P_\sigma, P_\tau$ are isomorphic iff there exists a bijection
preserving edge direction and transformation types.
\end{definition}

\begin{theorem}[Replay Equivalence]
Isomorphic provenance graphs reconstruct identical semantic states.
\end{theorem}

\subsection{Consequences}

\begin{corollary}[No Silent Rewrite]
Every semantic change is explicitly represented in the graph.
\end{corollary}

\begin{corollary}[Accountable Evolution]
All semantic states admit finite causal explanation.
\end{corollary}

\section{Semantic Neighborhood Topology}

\subsection{Latent Semantic Space}

Let $(X,d)$ be a complete, separable metric space (Polish space), called the
\emph{semantic latent space}.  

Let
\[
\phi : \Sigma \to X
\]
be an embedding mapping semantic objects to points in $X$.

\subsection{Embedding Invariants}

\begin{definition}[Embedding Consistency]
An embedding $\phi$ is consistent iff:
\[
\sigma_1 = \sigma_2 \;\Rightarrow\; \phi(\sigma_1) = \phi(\sigma_2).
\]
\end{definition}

\begin{definition}[Provenance Sensitivity]
$\phi$ is provenance-sensitive iff:
\[
P_{\sigma_1} \neq P_{\sigma_2} \;\Rightarrow\; d(\phi(\sigma_1), \phi(\sigma_2)) > 0.
\]
\end{definition}

\subsection{Neighborhood Definition}

\begin{definition}[Semantic Neighborhood]
Given $\sigma \in \Sigma$ and radius $\epsilon > 0$, define:
\[
\mathcal{N}_\epsilon(\sigma)
=
\{ \tau \in \Sigma \mid d(\phi(\sigma), \phi(\tau)) \le \epsilon \}.
\]
\end{definition}

\begin{definition}[Local Density]
The local density at $\sigma$ is:
\[
\rho(\sigma,\epsilon) = |\mathcal{N}_\epsilon(\sigma)|.
\]
\end{definition}

\subsection{Locality Constraint}

\begin{definition}[Local Interaction Constraint]
All admissible transformations $T$ satisfy:
\[
d(\phi(\sigma), \phi(T(\sigma))) \le \delta_T
\]
for some transformation-specific constant $\delta_T$.
\end{definition}

\begin{theorem}[Locality of Semantic Evolution]
Semantic evolution is locally bounded in $(X,d)$.
\end{theorem}

\begin{proof}
Transformations violating locality would require unbounded semantic reinterpretation,
which exceeds admissible entropy budgets.
\end{proof}

\subsection{Navigation Operator}

\begin{definition}[Navigation Operator]
Define navigation as a partial function:
\[
\mathcal{N} : \Sigma \times \mathbb{R}^+ \to \mathcal{P}(\Sigma)
\]
given by neighborhood queries.
\end{definition}

\begin{definition}[Traversal Path]
A traversal path is a sequence $(\sigma_0,\sigma_1,\dots,\sigma_n)$ such that:
\[
\sigma_{k+1} \in \mathcal{N}_\epsilon(\sigma_k).
\]
\end{definition}

\subsection{Spatial Memory}

\begin{definition}[Spatial Reaccessibility]
A semantic object $\sigma$ is spatially reaccessible iff:
\[
\forall t_1,t_2,\quad \phi_{t_1}(\sigma) = \phi_{t_2}(\sigma).
\]
\end{definition}

\begin{theorem}[Location Persistence]
Semantic objects occupy stable locations in $X$ across time.
\end{theorem}

\begin{proof}
Identity invariance and entropy-bounded transformations prevent embedding drift.
\end{proof}

\subsection{Feed Impossibility}

\begin{definition}[Global Feed Ordering]
A feed is a total order:
\[
\prec_F \; \subset \; \Sigma \times \Sigma
\]
intended to represent relevance.
\end{definition}

\begin{theorem}[No Global Feed Theorem]
There exists no total order $\prec_F$ compatible with all neighborhood relations.
\end{theorem}

\begin{proof}
Metric locality induces incompatible partial orders across disjoint neighborhoods.
\end{proof}

\subsection{Attention Conservation}

\begin{definition}[Attention Budget]
Let $A_u$ denote a finite attention budget for user $u$.
\end{definition}

\begin{proposition}[Local Attention Allocation]
Attention expenditure satisfies:
\[
\sum_{\sigma \in \mathcal{N}_\epsilon(\sigma_u)} a(\sigma) \le A_u.
\]
\end{proposition}

This prohibits unbounded attention capture.

\subsection{Dimensional Extension}

\begin{definition}[Higher-Dimensional Navigation]
Let $X = \prod_{k=1}^n X_k$ be a product space. Navigation may occur along projections
$\pi_k : X \to X_k$.
\end{definition}

\begin{theorem}[4× Interface Compatibility]
All navigation operators commute with dimensional projections:
\[
\pi_k(\mathcal{N}_\epsilon(\sigma)) = \mathcal{N}_{\epsilon_k}(\pi_k(\sigma)).
\]
\end{theorem}

\subsection{Consequences}

\begin{corollary}[No Infinite Scroll]
There exists no path $(\sigma_n)$ with strictly increasing novelty and bounded locality.
\end{corollary}

\begin{corollary}[Spatial Sociality]
Social interaction is equivalent to co-navigation in semantic space.
\end{corollary}

\section{Trust Accumulation Model}

\subsection{Preliminaries}

Let $\Sigma$ be the semantic state space and let $\mathcal{U}$ denote the set of identities
(users or agents). Each semantic object $\sigma \in \Sigma$ is associated with an identity
$i(\sigma) \in \mathcal{U}$.

Let $t \in \mathbb{R}_{\ge 0}$ denote semantic time as induced by provenance order.

\subsection{Persistence Measure}

\begin{definition}[Accessibility Indicator]
Define the accessibility indicator
\[
\chi_\sigma(t) =
\begin{cases}
1 & \text{if } \sigma \text{ is accessible at time } t \\
0 & \text{otherwise}.
\end{cases}
\]
\end{definition}

\begin{definition}[Persistence Functional]
Define persistence of a semantic object as
\[
\pi(\sigma) := \int_{0}^{\infty} \chi_\sigma(t)\, dt.
\]
\end{definition}

Persistence measures duration of availability rather than visibility.

\subsection{Entropy Stability}

\begin{definition}[Entropy Trajectory]
For a semantic object $\sigma$, define its entropy trajectory
\[
E_\sigma(t) := E(\sigma_t),
\]
where $\sigma_t$ denotes the semantic state at time $t$.
\end{definition}

\begin{definition}[Entropy Stability Functional]
Define entropy stability as
\[
S_E(\sigma) := -\int_{0}^{\infty} \left| \frac{d}{dt} E_\sigma(t) \right| dt.
\]
\end{definition}

Higher values of $S_E$ correspond to lower volatility.

\subsection{Reconciliation Credit}

\begin{definition}[Reconciliation Event]
A reconciliation event is a merge transformation
\[
M : (\sigma_1, \sigma_2) \mapsto \sigma_m
\]
that reduces entropy relative to the forked branches.
\end{definition}

\begin{definition}[Reconciliation Count]
Let $R(\sigma)$ denote the number of reconciliation events in $P_\sigma$.
\end{definition}

\subsection{Trust Functional}

\begin{definition}[Trust Functional]
Define trust of a semantic object as
\[
T(\sigma) := F\bigl(\pi(\sigma), S_E(\sigma), R(\sigma)\bigr),
\]
where $F : \mathbb{R}_{\ge 0}^3 \to \mathbb{R}_{\ge 0}$ is monotone in all arguments.
\end{definition}

Trust is defined structurally, independent of audience size.

\subsection{Identity-Level Trust}

\begin{definition}[Identity Trust]
Define trust of an identity $u \in \mathcal{U}$ as
\[
T(u) := \sum_{\sigma \in \Sigma : i(\sigma)=u} w(\sigma)\, T(\sigma),
\]
where $w(\sigma)$ is a normalization weight.
\end{definition}

\subsection{Popularity Independence}

\begin{definition}[Popularity Signals]
Let $L,S,F \in \mathbb{R}_{\ge 0}$ denote likes, shares, and followers respectively.
\end{definition}

\begin{theorem}[Popularity Irrelevance]
\[
\frac{\partial T(\sigma)}{\partial L}
=
\frac{\partial T(\sigma)}{\partial S}
=
\frac{\partial T(\sigma)}{\partial F}
=
0.
\]
\end{theorem}

\begin{proof}
Popularity variables do not appear in the definition of $T$.
\end{proof}

\subsection{Monotonicity Properties}

\begin{theorem}[Trust Monotonicity Under Valid Evolution]
If $\sigma \prec \sigma'$ and
\[
E(\sigma') \le E(\sigma),
\]
then
\[
T(\sigma') \ge T(\sigma).
\]
\end{theorem}

\begin{proof}
Persistence, entropy stability, and reconciliation count are non-decreasing under valid
transformations.
\end{proof}

---

\subsection{Adversarial Resistance}

\begin{definition}[Trust Inflation Attempt]
A trust inflation attempt is a transformation sequence increasing $T$ without increasing
$\pi$, $S_E$, or $R$.
\end{definition}

\begin{theorem}[No Trust Inflation]
Trust inflation attempts are impossible.
\end{theorem}

\begin{proof}
All arguments of $F$ correspond to structurally constrained quantities. None can be
artificially increased without satisfying invariants.
\end{proof}

\subsection{Temporal Convergence}

\begin{theorem}[Trust Convergence]
For any semantic object $\sigma$ with finite entropy budget,
\[
\lim_{t \to \infty} T(\sigma_t) \text{ exists}.
\]
\end{theorem}

\begin{proof}
Each component of $T$ is monotone and bounded.
\end{proof}

\subsection{Consequences}

\begin{corollary}[No Follower Economy]
There exists no mechanism by which follower accumulation increases trust.
\end{corollary}

\begin{corollary}[Time-Productive Interaction]
Time spent producing coherent semantic evolution strictly dominates time spent monitoring
visibility metrics.
\end{corollary}

\section{Spherepop Calculus Interface}

\subsection{Syntactic Categories}

Let the Spherepop language be defined by the following syntactic categories:
\[
\begin{aligned}
\sigma &\in \Sigma && \text{semantic objects} \\
m &\in \mathcal{M} && \text{modalities} \\
p &\in \mathcal{P} && \text{programs} \\
T &\in \mathcal{T} && \text{primitive transformations}
\end{aligned}
\]

Programs are defined inductively as:
\[
p ::= \text{id} \mid T \mid p_1 ; p_2 \mid p_1 \oplus p_2
\]

\subsection{Type System}

\begin{definition}[Typing Context]
A typing context is a tuple
\[
\Gamma = (\mathcal{M}_{\text{in}}, \mathcal{M}_{\text{out}}, \varepsilon)
\]
representing required input modalities, produced output modalities, and entropy budget.
\end{definition}

\begin{definition}[Typing Judgment]
Typing judgments have the form
\[
\Gamma \vdash p : \Sigma \to \Sigma.
\]
\end{definition}

\subsection{Primitive Typing Rules}

\[
\frac{}{\Gamma \vdash \text{id} : \Sigma \to \Sigma}
\quad
\text{(T-ID)}
\]

\[
\frac{
T \in \mathcal{T}
\quad
\mathcal{M}_{\text{in}}(T) \subseteq \mathcal{M}_\sigma
\quad
\varepsilon_T \le \varepsilon
}{
\Gamma \vdash T : \Sigma \to \Sigma
}
\quad
\text{(T-PRIM)}
\]

\subsection{Composition Rules}

\[
\frac{
\Gamma_1 \vdash p_1 : \Sigma \to \Sigma
\quad
\Gamma_2 \vdash p_2 : \Sigma \to \Sigma
}{
\Gamma_1 \cup \Gamma_2 \vdash p_1 ; p_2 : \Sigma \to \Sigma
}
\quad
\text{(T-SEQ)}
\]

\[
\frac{
\Gamma_1 \vdash p_1 : \Sigma \to \Sigma
\quad
\Gamma_2 \vdash p_2 : \Sigma \to \Sigma
}{
\max(\Gamma_1,\Gamma_2) \vdash p_1 \oplus p_2 : \Sigma \to \Sigma
}
\quad
\text{(T-ALT)}
\]

\subsection{Operational Semantics}

Define a small-step operational semantics:
\[
\langle p, \sigma \rangle \to \sigma'
\]

\[
\langle \text{id}, \sigma \rangle \to \sigma
\]

\[
\langle T, \sigma \rangle \to T(\sigma)
\quad \text{if } E(T(\sigma)) \le E(\sigma) + \varepsilon_T
\]

\[
\langle p_1 ; p_2, \sigma \rangle \to \langle p_2, \sigma' \rangle
\quad \text{if } \langle p_1, \sigma \rangle \to \sigma'
\]

\[
\langle p_1 \oplus p_2, \sigma \rangle \to \sigma'
\quad \text{if either branch yields valid output}
\]

\subsection{Entropy Preservation}

\begin{lemma}[Stepwise Entropy Bound]
If
\[
\langle p, \sigma \rangle \to \sigma'
\]
and
\[
\Gamma \vdash p : \Sigma \to \Sigma,
\]
then
\[
E(\sigma') \le E(\sigma) + \varepsilon.
\]
\end{lemma}

\subsection{Type Preservation}

\begin{theorem}[Subject Reduction]
If
\[
\Gamma \vdash p : \Sigma \to \Sigma
\quad \text{and} \quad
\langle p, \sigma \rangle \to \sigma',
\]
then $\sigma' \in \Sigma$.
\end{theorem}

\begin{proof}
By induction on the structure of $p$ using primitive and composition rules.
\end{proof}

\subsection{Progress}

\begin{theorem}[Progress]
If $\Gamma \vdash p : \Sigma \to \Sigma$ and $\sigma \in \Sigma$, then either
\[
p = \text{id}
\quad \text{or} \quad
\exists \sigma'.\; \langle p, \sigma \rangle \to \sigma'.
\]
\end{theorem}

\subsection{Confluence Properties}

\begin{definition}[Confluent Program]
A program $p$ is confluent iff all valid evaluation paths produce extensionally equal
semantic objects.
\end{definition}

\begin{theorem}[Local Confluence Under Entropy Bounds]
If all primitive transformations in $p$ commute and respect entropy bounds, then $p$ is
confluent.
\end{theorem}

\subsection{Consequences}

\begin{corollary}[No Hidden Side Effects]
Spherepop programs cannot modify semantic state outside explicit transformations.
\end{corollary}

\begin{corollary}[Programmable Social Interaction]
All social interaction reduces to well-typed program execution over $\Sigma$.
\end{corollary}

\section{Deterministic Substrate (Spherepop OS)}

\subsection{Event Log}

Let $\mathcal{E}$ be a countable set of events.

\begin{definition}[Event]
An event is a tuple
\[
e = (t, u, p, \sigma_{\text{in}}, \sigma_{\text{out}})
\]
where:
\begin{align*}
t &\in \mathbb{R}_{\ge 0} \quad \text{(timestamp)} \\
u &\in \mathcal{U} \quad \text{(identity)} \\
p &\in \mathcal{P} \quad \text{(Spherepop program)} \\
\sigma_{\text{in}}, \sigma_{\text{out}} &\in \Sigma.
\end{align*}
\end{definition}

\begin{definition}[Event Log]
An event log is a sequence
\[
\mathcal{L} = \langle e_1, e_2, \dots \rangle
\]
totally ordered by timestamp.
\end{definition}

\subsection{Kernel State}

\begin{definition}[Kernel State]
A kernel state is a tuple
\[
K = (O, U, R, M)
\]
where:
\begin{align*}
O &\subset \Sigma \quad \text{(active semantic objects)} \\
U &\subset \mathcal{U} \quad \text{(identities)} \\
R &\subset O \times O \quad \text{(references)} \\
M &\subset \mathcal{M} \quad \text{(active modalities)}.
\end{align*}
\end{definition}

\subsection{State Transition Function}

\begin{definition}[Transition Function]
Define the deterministic transition function
\[
\delta : K \times \mathcal{E} \to K
\]
such that:
\[
\delta(K, e) = K'
\]
where $K'$ is obtained by validating and applying $p(e)$ to $\sigma_{\text{in}}$.
\end{definition}

\begin{definition}[Validity Condition]
An event $e$ is valid iff:
\[
\Gamma \vdash p(e) : \Sigma \to \Sigma
\quad \text{and} \quad
\sigma_{\text{out}} = p(e)(\sigma_{\text{in}}).
\]
\end{definition}

\subsection{Replay Semantics}

\begin{definition}[Replay Operator]
Define replay as:
\[
\text{Replay}(\mathcal{L}_{\le n}) = \delta(\dots \delta(\delta(K_0, e_1), e_2), \dots, e_n)
\]
for fixed initial state $K_0$.
\end{definition}

\begin{theorem}[Replay Determinism]
For any prefix $\mathcal{L}_{\le n}$, the resulting kernel state is unique.
\end{theorem}

\begin{proof}
$\delta$ is deterministic and events are totally ordered.
\end{proof}

\subsection{Late-Joiner Equivalence}

\begin{definition}[Late Joiner]
A late joiner is an observer reconstructing state from $\mathcal{L}_{\le n}$.
\end{definition}

\begin{theorem}[Late-Joiner Equivalence]
All observers replaying $\mathcal{L}_{\le n}$ reconstruct identical kernel states.
\end{theorem}

\begin{proof}
Immediate from replay determinism.
\end{proof}

\subsection{Fork Semantics}

\begin{definition}[Explicit Fork]
A fork occurs only if two events reference the same $\sigma_{\text{in}}$ with incompatible
outputs.
\end{definition}

\begin{theorem}[No Implicit Forking]
Forks cannot arise without explicit event divergence in $\mathcal{L}$.
\end{theorem}

\begin{proof}
Sequential replay enforces a single outcome unless separate events encode divergence.
\end{proof}

\subsection{Causality and Order Independence}

\begin{lemma}[Local Commutativity]
Two events $e_i, e_j$ commute iff:
\[
\sigma_{\text{out}}(e_i) \cap \sigma_{\text{out}}(e_j) = \varnothing.
\]
\end{lemma}

\begin{theorem}[Causal Consistency]
Reordering commuting events yields equivalent kernel states.
\end{theorem}

\subsection{Impossibility Results}

\begin{theorem}[No State Rewrite]
There exists no event sequence that rewrites prior kernel state without leaving a trace in
$\mathcal{L}$.
\end{theorem}

\begin{theorem}[No Hidden Authority]
No identity can alter kernel state without producing a valid event.
\end{theorem}

\subsection{Consequences}

\begin{corollary}[Auditability]
All semantic state changes admit finite audit trails.
\end{corollary}

\begin{corollary}[Deterministic Social Memory]
Social state is fully reconstructible from the event log.
\end{corollary}

\section{Computational Complexity}

\subsection{Parameters}

Let:
\begin{align*}
n &= |\Sigma| \quad \text{(number of semantic objects)} \\
k &= \max_{\sigma} |P_\sigma| \quad \text{(max provenance size)} \\
\rho &= \max_{\sigma,\epsilon} \rho(\sigma,\epsilon) \quad \text{(max local density)} \\
m &= |\mathcal{L}| \quad \text{(event log length)} \\
u &= |\mathcal{U}| \quad \text{(identities)}.
\end{align*}

\subsection{Transformation Validation}

\begin{theorem}[Transformation Validation Complexity]
Validating a semantic transformation has time complexity
\[
O(k)
\]
and space complexity
\[
O(1).
\]
\end{theorem}

\begin{proof}
Validation requires checking modality completeness and provenance monotonicity, each linear
in provenance size.
\end{proof}


\subsection{Entropy Evaluation}

\begin{theorem}[Entropy Evaluation Complexity]
Computing $E(\sigma)$ has complexity
\[
O(|C_\sigma|).
\]
\end{theorem}

\begin{corollary}[Incremental Entropy Update]
Under bounded transformations,
\[
\Delta E = O(1).
\]
\end{corollary}

\subsection{Neighborhood Queries}

\begin{definition}[Neighborhood Query]
A neighborhood query returns
\[
\mathcal{N}_\epsilon(\sigma).
\]
\end{definition}

\begin{theorem}[Neighborhood Query Complexity]
Using approximate nearest neighbor indexing, neighborhood queries execute in
\[
O(\log n + \rho).
\]
\end{theorem}

\subsection{Navigation Paths}

\begin{definition}[Traversal Length]
Let $\ell$ be the length of a navigation path.
\end{definition}

\begin{proposition}[Traversal Cost]
Navigation cost is
\[
O(\ell(\log n + \rho)).
\]
\end{proposition}

---

\subsection{Replay and Reconstruction}

\begin{theorem}[Replay Complexity]
Replaying an event log prefix of length $m$ has time complexity
\[
O(mk).
\]
\end{theorem}

\begin{corollary}[Checkpointed Replay]
With periodic snapshots, replay complexity reduces to
\[
O((m - m_0)k)
\]
for nearest checkpoint $m_0$.
\end{corollary}

\subsection{Feed-Based Ranking Comparison}

\begin{definition}[Feed Ranking]
A feed ranking step computes a total order over candidate items.
\end{definition}

\begin{theorem}[Feed Ranking Complexity]
Feed ranking requires
\[
\Omega(n \log n)
\]
per recomputation.
\end{theorem}

\begin{theorem}[Global Recompute Impossibility]
Continuous feed ranking yields superlinear system-wide cost as $n \to \infty$.
\end{theorem}

\subsection{Locality Scaling Law}

\begin{theorem}[Locality Scaling]
Total system load scales as
\[
O(u \cdot \rho)
\]
independent of $n$.
\end{theorem}

\begin{proof}
Users interact only within bounded neighborhoods.
\end{proof}

\subsection{Memory Complexity}

\begin{theorem}[Structural Sharing Bound]
Total storage cost is
\[
O(n + mk)
\]
with persistent structure sharing.
\end{theorem}

\subsection{Consequences}

\begin{corollary}[No Global Bottlenecks]
No operation requires full-system traversal.
\end{corollary}

\begin{corollary}[Asymptotic Stability]
System performance remains stable under unbounded growth.
\end{corollary}

\section{Structural Impossibility Results}

\subsection{Preliminaries}

Assume the semantic system $(\Sigma, P, E, \phi)$ defined in Appendices A–H, with:
\[
\Sigma \text{ semantic objects}, \quad
P \text{ provenance DAGs}, \quad
E \text{ entropy functional}, \quad
\phi : \Sigma \to X.
\]

\subsection{Impossibility of Global Feeds}

\begin{definition}[Global Feed]
A global feed is a total order
\[
\prec_F \subset \Sigma \times \Sigma
\]
intended to represent relevance across all semantic objects.
\end{definition}

\begin{theorem}[Feed Impossibility]
There exists no $\prec_F$ such that:
\[
\sigma_1 \prec_F \sigma_2 \Rightarrow
d(\phi(\sigma_1), \phi(\sigma_2)) \le \epsilon
\quad \forall \sigma_1,\sigma_2.
\]
\end{theorem}

\begin{proof}
Total orders require comparability across arbitrarily distant neighborhoods, violating
locality constraints in $X$.
\end{proof}

\subsection{Impossibility of Viral Amplification}

\begin{definition}[Viral Amplification]
Viral amplification is the process of increasing visibility without increasing semantic
structure.
\end{definition}

\begin{theorem}[No Viral Amplification]
There exists no admissible transformation sequence $(T_k)$ such that:
\[
E(T_n(\sigma)) \le E(\sigma)
\quad \text{and} \quad
\text{visibility}(T_n(\sigma)) \to \infty.
\]
\end{theorem}

\begin{proof}
Visibility increase without structural information violates entropy bounds.
\end{proof}

\subsection{Impossibility of Reputation Laundering}

\begin{definition}[Reputation Laundering]
Reputation laundering is the increase of trust $T(u)$ without corresponding semantic
persistence.
\end{definition}

\begin{theorem}[No Reputation Laundering]
There exists no sequence of admissible transformations increasing $T(u)$ while keeping
$\pi(\sigma)$ bounded.
\end{theorem}

\begin{proof}
Trust depends monotonically on persistence and entropy stability.
\end{proof}

\subsection{Impossibility of Hidden Moderation}

\begin{definition}[Hidden Moderation]
Hidden moderation is the alteration or suppression of semantic objects without explicit
provenance.
\end{definition}

\begin{theorem}[No Hidden Moderation]
Hidden moderation is impossible.
\end{theorem}

\begin{proof}
All state transitions require explicit events in the log.
\end{proof}

\subsection{Impossibility of Shadow Banning}

\begin{definition}[Shadow Banning]
Shadow banning is selective visibility reduction without semantic transformation.
\end{definition}

\begin{theorem}[No Shadow Banning]
There exists no admissible operator that selectively reduces accessibility without
altering provenance.
\end{theorem}

\begin{proof}
Accessibility is a function of spatial position, not ranking.
\end{proof}

\subsection{Impossibility of Engagement Optimization}

\begin{definition}[Engagement Metric]
An engagement metric is a function
\[
g : \Sigma \to \mathbb{R}
\]
optimized independently of semantic structure.
\end{definition}

\begin{theorem}[No Engagement Objective]
There exists no engagement objective compatible with entropy-bounded evolution.
\end{theorem}

\begin{proof}
Engagement maximization incentivizes entropy increase.
\end{proof}

\subsection{Impossibility of Attention Extraction}

\begin{definition}[Attention Extraction]
Attention extraction is the unbounded consumption of user attention without semantic gain.
\end{definition}

\begin{theorem}[Attention Conservation]
For any user $u$,
\[
\int_0^\infty A_u(t)\, dt < \infty.
\]
\end{theorem}

\subsection{Impossibility of Meaning Collapse}

\begin{definition}[Meaning Collapse]
Meaning collapse is the convergence of semantic objects toward indistinguishable content.
\end{definition}

\begin{theorem}[No Meaning Collapse]
Entropy-bounded evolution prevents semantic homogenization.
\end{theorem}

\subsection{Separation Theorem}

\begin{theorem}[Architectural Separation]
The coherence--first semantic architecture is not a refinement of feed-based platforms, but
is structurally disjoint from them.
\end{theorem}

\begin{proof}
Feed-based systems require global ranking and engagement optimization, both of which are
impossible under the constraints established above.
\end{proof}

\subsection{Final Corollaries}

\begin{corollary}[No Ad-Based Incentive Compatibility]
Advertising optimization is incompatible with semantic coherence.
\end{corollary}

\begin{corollary}[Time Recovery]
User time expenditure produces persistent semantic value.
\end{corollary}

\begin{thebibliography}{99}

\bibitem{shannon1948}
C. E. Shannon.
\newblock A Mathematical Theory of Communication.
\newblock \emph{Bell System Technical Journal}, 27(3):379--423, 1948.

\bibitem{kolmogorov1965}
A. N. Kolmogorov.
\newblock Three Approaches to the Quantitative Definition of Information.
\newblock \emph{Problems of Information Transmission}, 1(1):1--7, 1965.

\bibitem{csikszentmihalyi1990}
M. Csikszentmihalyi.
\newblock \emph{Flow: The Psychology of Optimal Experience}.
\newblock Harper \& Row, 1990.

\bibitem{boyd2010}
d. boyd.
\newblock Social Network Sites as Networked Publics.
\newblock In Z. Papacharissi (ed.), \emph{A Networked Self}.
\newblock Routledge, 2010.

\bibitem{lampe2014}
C. Lampe, R. Wash, A. Velasquez, and E. Ozkaya.
\newblock Motivations to Participate in Online Communities.
\newblock \emph{CHI}, 2014.

\bibitem{grudin2010}
J. Grudin.
\newblock Groupware and Social Dynamics.
\newblock \emph{Communications of the ACM}, 53(5):66--74, 2010.

\bibitem{viegas2007}
F. Vi{\'e}gas, M. Wattenberg.
\newblock Visualization of Online Social Interactions.
\newblock \emph{CHI}, 2007.

\bibitem{thompson1967}
K. Thompson.
\newblock Reflections on Trusting Trust.
\newblock \emph{Communications of the ACM}, 27(8):761--763, 1984.

\bibitem{git2005}
L. Torvalds, J. Hamano.
\newblock Git: Fast Version Control System.
\newblock 2005. \url{https://git-scm.com}

\bibitem{raymond1999}
E. S. Raymond.
\newblock \emph{The Cathedral and the Bazaar}.
\newblock O’Reilly Media, 1999.

\bibitem{bernerslee2001}
T. Berners-Lee, J. Hendler, O. Lassila.
\newblock The Semantic Web.
\newblock \emph{Scientific American}, 284(5):28--37, 2001.

\bibitem{halpin2010}
H. Halpin, P. J. Hayes.
\newblock When Ontologies Collide.
\newblock \emph{Journal of Web Semantics}, 8(2--3):113--128, 2010.

\bibitem{luhmann1995}
N. Luhmann.
\newblock \emph{Social Systems}.
\newblock Stanford University Press, 1995.

\bibitem{latour2005}
B. Latour.
\newblock \emph{Reassembling the Social}.
\newblock Oxford University Press, 2005.

\bibitem{ostrom1990}
E. Ostrom.
\newblock \emph{Governing the Commons}.
\newblock Cambridge University Press, 1990.

\bibitem{floridi2011}
L. Floridi.
\newblock \emph{The Philosophy of Information}.
\newblock Oxford University Press, 2011.

\bibitem{bowker2000}
G. C. Bowker, S. L. Star.
\newblock \emph{Sorting Things Out: Classification and Its Consequences}.
\newblock MIT Press, 2000.

\bibitem{manovich2001}
L. Manovich.
\newblock \emph{The Language of New Media}.
\newblock MIT Press, 2001.

\bibitem{kleinberg1999}
J. Kleinberg.
\newblock Authoritative Sources in a Hyperlinked Environment.
\newblock \emph{Journal of the ACM}, 46(5):604--632, 1999.

\bibitem{page1999}
L. Page, S. Brin.
\newblock The Anatomy of a Large-Scale Hypertextual Web Search Engine.
\newblock \emph{WWW}, 1998.

\bibitem{pariser2011}
E. Pariser.
\newblock \emph{The Filter Bubble}.
\newblock Penguin Press, 2011.

\bibitem{zuckerberg2017}
M. Zuckerberg.
\newblock Building Global Community.
\newblock Facebook Note, 2017.

\bibitem{foucault1977}
M. Foucault.
\newblock \emph{Discipline and Punish}.
\newblock Pantheon Books, 1977.

\bibitem{wiener1948}
N. Wiener.
\newblock \emph{Cybernetics}.
\newblock MIT Press, 1948.

\bibitem{ashby1956}
W. R. Ashby.
\newblock \emph{An Introduction to Cybernetics}.
\newblock Chapman \& Hall, 1956.

\bibitem{friston2010}
K. Friston.
\newblock The Free-Energy Principle.
\newblock \emph{Nature Reviews Neuroscience}, 11:127--138, 2010.

\bibitem{benkler2006}
Y. Benkler.
\newblock \emph{The Wealth of Networks}.
\newblock Yale University Press, 2006.

\bibitem{doctorow2023}
C. Doctorow.
\newblock The Enshittification of TikTok.
\newblock \emph{Pluralistic}, 2023.
.

\end{thebibliography}

\end{document}
